\section{Related Work}

\subsection{ECMO Treatment Prediction} 
To date, there exist a substantial gap between existing studies and ECMO treatment analysis. Most studies are limited to ECMO mortality scoring systems \cite{pappalardo2013predicting, schmidt2014predicting, schmidt2013preserve, shah2023validation} that rely on identifying pre-ECMO variables, which are available and validated only from patients already supported on ECMO. As such, none of these scores have utilized an appropriate matched non-ECMO cohort, rendering them incapable of identifying the  patients who should receive ECMO support \cite{shah2021extracorporeal,short2022extracorporeal}. A recent study uses gradient boosting trees to predict ECMO treatment assignment \cite{xue2022ecmo}, but it is not designed to predict treatment effect in terms of factual and counterfactual outcomes, which is an important contribution of this work. 

\subsection{Individual Treatment Effect} 
To predict each individual's treatment assignment and treatment effects, the most intuitive approach is to build separate single-task predictors for propensity (probability of receiving ECMO treatment), treatment outcome and control outcome (i.e., survival or death with/without ECMO treatment), respectively  \cite{alaa2017dcnpd,johansson2016bnn,shalit2017cfrnet,chipman2010bart}. Such models are generalized as meta-learners in recent literature \cite{curth2021meta, curth2021st}. These models are prone to the strong selection bias and label imbalance in ECMO treatment assignment, as each treatment-outcome predictor is only exposed to a specific patient group, but not the whole population.

 Another popular approach reported in the literature is by adapting non-parametric models for individual level treatment effect. The simplest version is the k nearest neighbors model, and more advanced models are adapted from tree-based ensemble models. Causal Forest has been developed from random forest to obtain a consistent estimator with semi-parametric asymptotic convergence rate \cite{wager2018cf}. Bayesian Additive Regression Trees (BART)-based methods have been proposed to build the trees with the regularization prior and the backfitting Markov chain Monte Carlo (MCMC) algorithm \cite{chipman2010bart, hahn2020bayesian}. Compared to neural network-based solutions, it is hard for such models to build and regularize the patients representations, hence the curse of dimensionality in characterizing scarce treatment group remains a challenge. 
 
 \subsubsection{Representation Learning} 
 Our proposed TVAE is related to earlier works using deep representation learning for treatment effect estimation. To build a balanced representation between treatment and control groups, various strategies are adopted to force information sharing between the groups. An popular strategy is to remove the selection bias from the representations of treatment and control groups, hence the representations of both groups are similar. Such representation can be from shared layers (such as SNet \cite{curth2021st}, DCN-PD \cite{alaa2017dcnpd}, Dragonnet \cite{shi2019adapting}, TARNET \cite{shalit2017tarnet}, BNN \cite{johansson2016bnn}) or separate networks (e.g., TNet \cite{curth2021st}). In contrast, TVAE acknowledges the strong selection bias in ECMO data hence does not force the representations of two groups to be the same. Instead, TVAE \textit{disentangles} the representation of patients into different aspects (or latent dimensions), and extracts the biased aspect to a designated latent dimension. By doing so, the remaining latent dimensions are naturally "balanced" as both groups share the similar information in these aspects. Considering the rareness of ECMO treatment events, TVAE further avoids the potential over-fitting by maximizing the cross-group similarities in these remaining latent dimensions. 
 
 \subsubsection{Deep Generative Models} 
 Another direction to estimate treatment effects is through deep generative models. CEVAE \cite{louizos2017cevae} and IntactVAE \cite{wu2021intact} adapt the Variational Autoencoder (VAE) to transform the representation of all patients into a common latent space, and build auxiliary networks to predict factual and counterfactual outcomes. GANITE, on the other hand, learn the counterfactual
distributions instead of conditional expected values \cite{yoon2018ganite}. 
As a VAE-based framework, TVAE possesses the salient properties of the abovementioned works, but in a different fashion. TVAE has a re-organized latent space that encodes patients' characteristics by explicitly expressing the predicted distributions of treatment assignment as well as factual and counterfactual outcomes. This eliminates the extra complexity of adding auxiliary networks and its potential insufficient training (as each auxiliary network for treatment outcome is only trained by a subset of data).  Such re-organized latent space further regulates the counterfactual outcomes through the clustering effect as well as input reconstruction. Moreover, TVAE leverages its generative feature to tackle the data imbalance in ECMO prediction: it augments ECMO cases by upsampling from their latent distributions, hence achieving label balance in training iterations. 
